%% kapa: exam=1 round=36 year=2025 period=1 form=A subject=real_estate subject_ko=부동산학원론 seq=23 number=103
%% source: Q-Net 감정평가사 기출문제 (cst00309) artlSeq=5231525 file=raw/1008/5231525/2025년 제36회 감정평가사 1차시험 1교시 문제지 원본.pdf
%% conversion: kordoc 4.24.51 -> tools/gen1.py (auto); review=proofread-sonnet; flags=none
\begin{question}{103}
부동산마케팅활동에 관한 설명으로 옳지 않은 것은?
\begin{choices}
  \item 부동산마케팅은 소비자들이 원하는 필요와 욕구를 반영하여 시장을 세분화하고 이를 바탕으로 부동산의 제품화, 가격산정, 입지선정 및 촉진활동 등 마케팅 전략을 세워서 부동산을 매매하고 임대차하는 일련의 과정을 말한다.
  \item 시장세분화란 전체 시장을 일정한 기준에 의해 동질적인 세분시장으로 구분하는 과정을 말한다.
  \item 표적시장이란 마케팅 환경변화에 대응하여 경쟁사와의 관계에서 자사가 보유한 역량과 자원으로 최대한의 시장성과를 얻을 수 있는 최적의 시장을 말한다.
  \item 포지셔닝이란 경쟁우위 달성을 위해 경쟁제품과 다르게 인식되도록 마케팅믹스를 사용하여 고객의 마음속에 제품의 위치를 심어주는 과정을 말한다.
  \item AIDA 원리는 고객의 구매의사 결정단계를 심리적 발전단계에 맞춘 것으로 행동(Action), 관심(Interest), 욕망(Desire), 주목(Attention)의 순서를 거친다.
\end{choices}
\end{question}
\answer{5}
